\documentclass[11pt]{amsart}
\usepackage[T1]{fontenc}
\usepackage{lmodern,amsmath,amssymb,amsthm,mathtools,microtype,booktabs,array,longtable}
\usepackage[margin=1.06in]{geometry}
\usepackage{tikz}
\usepackage{xr-hyper}
\usepackage[hidelinks]{hyperref}
\hypersetup{pdftitle={Finite-Use Discrimination Geometry of Ordered Quantum Measurements},pdfauthor={Qian Qi}}
\emergencystretch=2em
\numberwithin{equation}{section}
\newtheorem{theorem}{Theorem}[section]
\newtheorem{lemma}[theorem]{Lemma}
\newtheorem{proposition}[theorem]{Proposition}
\newtheorem{corollary}[theorem]{Corollary}
\theoremstyle{definition}\newtheorem{definition}[theorem]{Definition}
\theoremstyle{remark}\newtheorem{remark}[theorem]{Remark}
\newcommand{\Q}{\mathbb Q}\newcommand{\R}{\mathbb R}\newcommand{\C}{\mathbb C}\newcommand{\N}{\mathbb N}\newcommand{\Z}{\mathbb Z}\newcommand{\E}{\mathbb E}
\newcommand{\ip}[2]{\langle #1,#2\rangle}\newcommand{\norm}[1]{\lVert #1\rVert}
\newcommand{\epsc}{\varepsilon_{\mathrm c}}
\DeclareMathOperator{\osc}{osc}\DeclareMathOperator{\spanop}{span}\DeclareMathOperator{\End}{End}\DeclareMathOperator{\diag}{diag}\DeclareMathOperator{\tr}{tr}\DeclareMathOperator{\rad}{rad}\DeclareMathOperator{\TV}{TV}
\title[Finite-use discrimination geometry]{Finite-Use Discrimination Geometry of Ordered Quantum Measurements}
\author{Qian Qi}
\date{6 October 2026}
\subjclass[2020]{Primary 81P15, 81P47; Secondary 94A17, 81P70}
\externaldocument{build/xref-supplement}[BINARY_SUPPLEMENT.pdf]
\begin{document}
\raggedbottom
\begin{abstract}
We study ordered quantum measurements with separate acquisition and
receiver resources.  For arbitrary finite two-call decision experiments,
we express the optimal unit-reset value as a concave-envelope problem
on density matrices with a common barycenter.  The optimum is attained
with at most $d_1^2$ receiver-instrument outcomes per first label and
quantum receiver dimension $d_1d_2$.  An exact dual gives majorization
and contact certificates and identifies the effect of receiver feedback.
Pure-atom envelopes give the classical optimum and necessary and
sufficient certificates of retained-receiver advantage.
These results require neither a symmetric ensemble nor a special prior.
We then solve a finite basis-measurement experiment in every dimension
with equal hypothesis priors.  Complete measure-and-reprepare adaptation,
unit-reset receiver storage and unrestricted two-call acquisition have
respective optimal successes $1/2+t^2(d-1)/(2d(d+1))$,
$1/2+t^2(d-1)/(2d^2)$ and $1/2+t^2/(2d)$.
A centered filtered-swap inequality bounds the entire reset class;
explicit strategies attain all three values.  The hierarchy persists
under full-rank deformation and small changes of channels, priors and
the shared latent law.  A quantitative equality estimate controls the
Gram marginals of near-optimal reset strategies.  We retain the complementary local geometric
classification: in a fixed quadratic tangent tube, an allocation with
$N=\sum n_r$ and $V=\sum n_r^2$ has trace orders
$\min\{1,\sqrt Ns\}$, $\min\{1,\sqrt Vs\}$ or $\min\{1,Ns\}$,
according to the complete support-covariance trichotomy.
Hard reset-policy square budget $Q$ yields coherent order
$\min\{1,\sqrt Qs\}$, uniformly over the allocation and unknown quadratic
remainder.  Local constants depend on the fixed tube; exact variational
values do not assert strict advantage for every fixed pair.
\end{abstract}
\maketitle
\input{editions/operational-introduction91}
\input{sections/67-coupled-boundary-covariance}
\input{sections/72-finite-product-prerequisite}
\input{sections/69-support-kernel}
\input{sections/70-finite-angle-orbits}
\input{sections/73-smooth-boundary-curves}
\input{sections/74-certified-control-stability}
\input{sections/75-one-sided-tangent-neighborhoods}
\input{sections/76-independent-probes-and-higher-jets}
\input{sections/77-nonempty-tangent-neighborhoods}
\input{sections/78-entanglement-width}
\input{sections/79-classical-feedback-reset-width}
\input{sections/80-allocation-profiles}
\input{sections/81-quadratic-budgets-and-memory}
\input{sections/83-reset-variational-principle}
\input{sections/82-operational-memory-hierarchy}
\input{sections/84-equal-prior-memory-hierarchy}
\input{editions/current-comparison91}
\input{editions/bibliography91}
\end{document}
